% s-337-p337-m002; original PDF 337, 338
\Aphorism{s-337-p337-m002}{386}
Zwei kosmische Augenblicke, beide vielleicht nur in einem und demselben Gehirn möglich, beide kurz, schwer, unendlich wie Mittagsträume\Dots{3} Der eine: schwindelndes Entsetzen vor der Unbegreiflichkeit der Welt, marternde Unbefriedigung jedes Verlangens nach Schönheit, Güte, Vernunft. Alle Bänder des Zusammenhanges scheinen gelöst zu flattern; wohin die Phantasie mit klammernden Händen greift, zerbröckelt alles in sinnlose Einzelheiten, in farblose wärmelose gestaltlose Atome, die in ihrer grauen Unbeschaffenheit doch nicht verständlicher werden. Ein gähnender Zweifel, der selbst die Identität schwanken macht; ist wirklich Gleiches noch gleich? Das andere Weltbild: tiefste, athemlose Ruhe und Sättigung, Dasein und Leben zu durchsichtiger Klarheit verglast, vergeistigt. Musik des Schweigens; wir würden jede irdische Musik durch sie hindurch hören, wie wir durch Glas sehen. Die Thatsache „Welt“ ist so undenkbar begreiflich, so unbegreiflich denkbar geworden, die Einheit aller Einheiten, und doch hat sie nichts von ihrer Vielfachheit verloren: jede Farbe leuchtet, jede Blüthe duftet, jedes Leiden brennt, jede Erlösung kühlt. Ist es die wissenschaftliche, die sittliche, die künstlerische, die religiöse Entschleierung des Welträthsels, die uns zu Theil geworden? Kaum, aber es scheint, als könnten wir alles, was wir wissen, zweifeln, lieben, anbeten, mit Einem Worte aussprechen. Mit welchem? — aber wenn wir darüber nachsinnen, ist der Traum bereits verflogen.
\par

